% Regression test for \morphName/\morphSeqName's auto-detection,
% mirroring test-catname-decoration.tex. Covers every distinct
% argument actually found in LCS.arXiv.V2.tex/M-atlas.tex's
% funcname/funcseqname call sites (11 unique arguments across 5409
% combined calls) -- including funcseqname{f^2}, the one real site a
% single-token-only test would have mis-classified (see
% semantic-markup.dtx, the "Tightened, 2026-09-01" correction under
% \morphName's documentation).
\documentclass{article}
\usepackage{semantic-markup}

\ExplSyntaxOn

\let\reservedmathit\mathit
\let\reservedmathrm\mathrm
\tl_new:N \l_lsmtest_got_tl
\renewcommand{\mathit}[1]{\tl_gset:Nn \l_lsmtest_got_tl {italic} \reservedmathit{#1}}
\renewcommand{\mathrm}[1]{\tl_gset:Nn \l_lsmtest_got_tl {roman} \reservedmathrm{#1}}
% \bm (from the bm package) does its own internal re-scanning of its
% argument that doesn't tolerate the side-effecting \mathit/\mathrm
% redefinitions above -- irrelevant to what this test actually checks
% (which branch \lsm_if_arbitrary_symbol:nTF took), so bypass it with
% a transparent pass-through for the duration of this test only.
\renewcommand{\bm}[1]{#1}

% #1 = macro to call (morphName or morphSeqName), #2 = argument,
% #3 = expected (italic|roman)
\cs_new_protected:Npn \lsmtestMorph #1#2#3
  {
    \tl_clear:N \l_lsmtest_got_tl
    \hbox{$#1{#2}$}
    \tl_if_eq:NnTF \l_lsmtest_got_tl {#3}
      {
        \iow_term:x
          { ok~~ \tl_to_str:n{#1} \{ \tl_to_str:n{#2} \}~~->~~#3 }
      }
      {
        \iow_term:x
          { FAIL~~ \tl_to_str:n{#1} \{ \tl_to_str:n{#2} \}~~expected~#3~~got~
            \l_lsmtest_got_tl }
        \msg_error:nnxx { lsmtest } { mismatch }
          { \tl_to_str:n{#2} } {#3}
      }
  }
\msg_new:nnn { lsmtest } { mismatch }
  { morph\{#1\}~expected~#2~but~did~not~get~it. }

\ExplSyntaxOff

\begin{document}

% Single-token arguments (called via funcname and/or funcseqname in
% the papers -- expect italic / bold italic), including the two
% single-control-sequence arguments \pi and \rho.
\lsmtestMorph{\morphName}{F}{italic}
\lsmtestMorph{\morphSeqName}{F}{italic}
\lsmtestMorph{\morphSeqName}{G}{italic}
\lsmtestMorph{\morphName}{\pi}{italic}
\lsmtestMorph{\morphSeqName}{\pi}{italic}
\lsmtestMorph{\morphSeqName}{\rho}{italic}
\lsmtestMorph{\morphName}{f}{italic}
\lsmtestMorph{\morphSeqName}{f}{italic}
\lsmtestMorph{\morphName}{g}{italic}
\lsmtestMorph{\morphSeqName}{g}{italic}
\lsmtestMorph{\morphName}{i}{italic}
\lsmtestMorph{\morphName}{m}{italic}
\lsmtestMorph{\morphSeqName}{m}{italic}
\lsmtestMorph{\morphName}{s}{italic}
\lsmtestMorph{\morphName}{t}{italic}

% The decorated symbol that a single-token-only test gets wrong
% (called via funcseqname despite not being a single token -- expect
% italic; this is the actual regression being guarded).
\lsmtestMorph{\morphSeqName}{f^2}{italic}

All 11 real-world arguments classified correctly (16 calls across
morphName/morphSeqName as actually used in the source papers).
\end{document}
